% !TEX program = xelatex
% ============================================================
%  用 LaTeX beamer 做 PPT —— 最小可运行模板
%  编译方式：必须使用 XeLaTeX
%      xelatex beamer_minimal_template.tex
%  或在 VS Code (LaTeX Workshop) / TeXstudio 中选择 XeLaTeX
% ============================================================
\documentclass[aspectratio=169, 11pt]{beamer}   % 16:9 比例；4:3 用 aspectratio=43

% ---------- 中文支持：必须加 ctex（底层用 xeCJK） ----------
\usepackage{ctex}

% ---------- 插图 ----------
\usepackage{graphicx}

% ---------- 主题与配色 ----------
\usetheme{Madrid}          % 换成 metropolis / Warsaw / Berlin 等
\usecolortheme{default}    % 换成 whale / dolphin / seahorse / beetle 等

% ---------- 标题信息 ----------
\title{用 Beamer 做 PPT 最简入门}
\subtitle{一个可直接编译的最小示例}
\author{Your Name}
\institute{Your Institute}
\date{\today}

\begin{document}

% ========== 标题页 ==========
\begin{frame}
    \titlepage
\end{frame}

% ========== 目录页：\section 会自动生成目录条目 ==========
\begin{frame}{目录}
    \tableofcontents
\end{frame}

% ========== 第 1 节 ==========
\section{基础用法}

\begin{frame}{列表与 \texttt{\textbackslash pause} 动画}
    \begin{itemize}
        \item 第一点：itemize 生成无序列表
        \item 第二点：enumerate 可生成有序列表
        \item 第三点：用 \texttt{\textbackslash item} 逐条书写
    \end{itemize}

    \pause   % 到这里暂停，下一页内容按需分步出现

    按过 \texttt{\textbackslash pause} 之后，这一句才会出现。
\end{frame}

\begin{frame}{居中与换行}
    \centering
    这一行居中显示。\\
    这里换行用双反斜杠 \texttt{\textbackslash\textbackslash}。
\end{frame}

% ========== 第 2 节 ==========
\section{进阶布局}

\begin{frame}{分栏 columns 与高亮块 block}
    \begin{columns}[T]           % [T] 顶端对齐
        \begin{column}{0.48\textwidth}
            \begin{block}{左栏：普通高亮块}
                \begin{itemize}
                    \item 要点一
                    \item 要点二
                \end{itemize}
            \end{block}
        \end{column}
        \begin{column}{0.48\textwidth}
            \begin{alertblock}{右栏：警告高亮块}
                alertblock / exampleblock / block 三种样式可选。
            \end{alertblock}
        \end{column}
    \end{columns}
\end{frame}

\begin{frame}{插入图片}
    \begin{figure}
        \centering
        % 把下面的 example.png 换成你自己的文件名，并取消注释即可使用：
        % \includegraphics[width=0.6\textwidth]{example.png}
        \fbox{\parbox[c][3cm][c]{0.6\textwidth}{\centering 这里放图片\\（取消注释上面的 include graphics 行）}}
        \caption{图片标题示例}
    \end{figure}
\end{frame}

% ========== 第 3 节 ==========
\section{常见踩坑}

\begin{frame}{四句话总结}
    \begin{itemize}
        \item 必须用 \textbf{XeLaTeX} 编译；
        \item 中文要加 \texttt{\textbackslash usepackage\{ctex\}}；
        \item 标题里的特殊字符要转义；
        \item 换行用 \texttt{\textbackslash\textbackslash}。
    \end{itemize}
\end{frame}

\end{document}
